\documentclass[11pt,a4paper]{article}
\usepackage[margin=1in]{geometry}
\usepackage{times}
\usepackage{latexsym}
\usepackage{graphicx}
\usepackage{booktabs}
\usepackage{amsmath}
\usepackage{amssymb}
\usepackage{multirow}
\usepackage{xcolor}
\usepackage{tabularx}
\usepackage{arydshln}
\usepackage[T1]{fontenc}
\usepackage[utf8]{inputenc}
\usepackage{microtype}
\usepackage{natbib}
\usepackage{hyperref}

\title{Why Doesn't Symbolic Consensus Improve LLM Reasoning?\\A Diagnostic Study}

\author{Anonymous}

\date{}

\begin{document}
\maketitle

% ============================================================
\begin{abstract}
Self-Consistency (SC) aggregates multiple reasoning traces from large language models via majority voting, discarding the logical structure within each trace.
Symbolic Invariant Consensus Aggregation (SICA) attempts to improve upon SC by extracting first-order logic constraints from traces and using weighted partial MAX-SAT to find globally consistent answers.
We conduct a systematic diagnostic study of SICA across 7 models from 3 families (Qwen, Mistral, LLaMA), 2 reasoning architectures (QwQ, DeepSeek-R1), and 4 benchmarks (FOLIO, ProofWriter, GSM8K).
On 804 expanded evaluation instances, SICA achieves only +0.3 percentage points over SC ($p=0.42$), failing to demonstrate statistically significant improvement.
Our analysis reveals three findings:
(1)~the multi-candidate rate $C_0$---a zero-cost metric computed from SC votes alone---reliably distinguishes instruction models ($C_0>20\%$) from reasoning models ($C_0<5\%$) across all families, characterizing \emph{when} models disagree rather than \emph{whether} symbolic methods help;
(2)~constraint-based aggregation fails because LLM-extracted constraints are predominantly shallow syntactic patterns (only 45\% semantically meaningful), and MAX-SAT satisfiability checking adds no value over simple counting ($\Delta=0.1$pp, $p=0.48$);
(3)~reasoning models exhibit structural self-consistency ($C_0\approx0$) regardless of temperature, a property arising from extended chain-of-thought training rather than model family or scale.
These negative results inform practitioners about the current limitations of symbolic aggregation and provide a diagnostic framework for evaluating future consensus methods.
\end{abstract}

% ============================================================
\section{Introduction}
\label{sec:intro}

Self-Consistency \cite{wang2023selfconsistency} has become the standard approach for aggregating multiple reasoning traces from large language models (LLMs): sample $K$ completions and return the majority-voted answer.
While effective, SC treats each reasoning chain as an opaque vote, discarding the logical structure within.
When votes are tied---a frequent occurrence at moderate $K$---SC falls back to arbitrary rules (random selection, alphabetical ordering) with no principled resolution.

Symbolic Invariant Consensus Aggregation (SICA) was proposed to address this limitation by extracting first-order logic constraints from reasoning traces and using weighted partial MAX-SAT solving to identify the most logically consistent answer.
The intuition is compelling: if we can extract the logical invariants that each reasoning chain implicitly relies upon, we can aggregate at the level of logical structure rather than surface-level votes.

In this paper, we report the results of a systematic evaluation of SICA that yields a negative finding: \textbf{symbolic constraint aggregation does not reliably improve over majority voting} across models, scales, and domains.
Rather than treating this as a simple failure, we conduct a diagnostic investigation into \emph{why} constraint-based consensus fails and \emph{what} we can learn from this failure.

Our contributions are:

\begin{enumerate}
\item \textbf{$C_0$ as a zero-cost diversity indicator.}
We show that the multi-candidate rate $C_0$---the fraction of problems where SC produces more than one candidate answer---reliably characterizes model behavior across 3 model families.
Instruction-tuned models exhibit $C_0>20\%$ universally, while reasoning models show $C_0<5\%$ regardless of temperature.
$C_0$ reflects \emph{when} models disagree, serving as a practical diagnostic for the applicability space of any consensus method.

\item \textbf{Negative result with failure mode analysis.}
We demonstrate that SICA's apparent gains on small samples ($n=60$) are tiebreak artifacts that vanish at scale ($n=804$, $\Delta=+0.3$pp, $p=0.42$).
We identify the root cause: LLM-extracted constraints are predominantly shallow patterns (45\% semantically meaningful, 35\% shallow matches, 20\% noise), and MAX-SAT solving over these constraints provides no advantage over simple counting.

\item \textbf{Taxonomy of model behavior under sampling.}
We provide a systematic characterization of how instruction-tuned vs.\ reasoning-trained models behave under temperature sampling, showing that reasoning training fundamentally collapses output diversity---same LLaMA-8B base yields $C_0=30\%$ with instruction tuning vs.\ $C_0=5\%$ with reasoning distillation.
\end{enumerate}

% ============================================================
\section{Background}
\label{sec:background}

\paragraph{Self-Consistency.}
Given a question $q$, SC samples $K$ reasoning traces $(r_1, \dots, r_K)$ from an LLM at temperature $T>0$, extracts an answer $a_i$ from each trace, and returns $\hat{a} = \arg\max_a \sum_{i=1}^{K} \mathbf{1}[a_i = a]$.
When multiple answers tie for the maximum count, SC uses an arbitrary tiebreaking rule (random, alphabetical, or first-seen).

\paragraph{SICA.}
SICA extends SC with a constraint extraction and satisfiability-checking phase.
For each trace $r_i$, a constraint extractor (an LLM prompted for logical analysis) produces a set of first-order logic constraints $\mathcal{C}_i$.
These constraints are aggregated across traces sharing the same answer into a weighted partial MAX-SAT instance, where the weight of each constraint reflects its frequency across traces.
A MAX-SAT solver (Z3) identifies the answer whose supporting constraint set achieves the highest satisfiability score.
SICA reduces to SC when all constraints are trivially satisfiable or when a single answer dominates (no tie to break).

\paragraph{Multi-candidate rate ($C_0$).}
We define $C_0$ as the fraction of problems in a dataset where SC produces more than one candidate answer among $K$ traces:
\begin{equation}
C_0 = \frac{1}{N} \sum_{j=1}^{N} \mathbf{1}\left[|\{a : \exists i, a_i^{(j)} = a\}| > 1\right]
\end{equation}
where $a_i^{(j)}$ is the answer extracted from the $i$-th trace for problem $j$.
$C_0$ is computed from SC outputs alone at zero additional cost.

% ============================================================
\section{Method}
\label{sec:method}

The SICA pipeline operates in three phases:

\paragraph{Phase 1: Trace generation.}
We sample $K$ reasoning traces from the target LLM at temperature $T$, identical to SC.

\paragraph{Phase 2: Constraint extraction.}
For each trace $r_i$, we prompt an LLM to extract a set of first-order logic constraints $\mathcal{C}_i = \{c_{i,1}, \dots, c_{i,m}\}$ that the reasoning chain implicitly relies upon.
Each constraint $c$ is a predicate over the answer space (e.g., ``the answer must satisfy property $P$'').

\paragraph{Phase 3: MAX-SAT aggregation.}
For each candidate answer $a$, we collect all constraints from traces supporting $a$ and formulate a weighted partial MAX-SAT instance.
Hard constraints encode answer-format requirements; soft constraints are the extracted logical predicates, weighted by their frequency across traces.
The answer $\hat{a}$ with the highest MAX-SAT satisfiability score is selected:
\begin{equation}
\hat{a} = \arg\max_a \textsc{MaxSAT\text{-}Score}(\mathcal{C}_a, \mathbf{w}_a)
\end{equation}
where $\mathcal{C}_a = \bigcup_{i: a_i=a} \mathcal{C}_i$ and $\mathbf{w}_a$ are the frequency-based weights.
When the MAX-SAT score is identical across candidates (common when constraints are trivially satisfiable), SICA falls back to vote count, then to SC's tiebreaking rule.

% ============================================================
\section{Experimental Setup}
\label{sec:setup}

\paragraph{Datasets.}
We evaluate on FOLIO \cite{han2022folio}, a first-order logic natural language inference benchmark (204 instances, 3-class: True/False/Unknown);
ProofWriter \cite{tafjord2021proofwriter} depth-5 subset (600 instances, binary: Proved/Disproved);
and a GSM8K \cite{cobbe2021gsm8k} pilot (30 instances, open-answer arithmetic).

\paragraph{Models.}
We evaluate 7 models from 3 families spanning two paradigms:

\noindent\emph{Instruction-tuned:}
Qwen2.5-7B-Instruct, Qwen2.5-14B-Instruct, Qwen2.5-72B-Instruct-AWQ (4-bit),
Mistral-7B-Instruct-v0.3, LLaMA-3.1-8B-Instruct.

\noindent\emph{Reasoning-trained:}
QwQ-32B, DeepSeek-R1-Distill-LLaMA-8B.

\paragraph{Configuration.}
Unless noted otherwise: $K=12$ traces, temperature $T=0.7$.
Constraint extraction uses the same model as trace generation.
MAX-SAT solving uses Z3.
All experiments use vLLM for efficient batched inference.

% ============================================================
\section{Results}
\label{sec:results}

\subsection{Main Results: SICA Does Not Improve Over SC at Scale}
\label{sec:main_results}

Table~\ref{tab:main} presents the primary evaluation on 804 instances (FOLIO + ProofWriter-D5) with Qwen2.5-14B-Instruct.
SICA achieves 72.3\% vs.\ SC's 72.0\%, a gap of +0.3 percentage points that is not statistically significant (bootstrap $p=0.42$).

Breaking down by dataset: on FOLIO (204 instances), SICA shows +1.5pp ($p=0.28$); on ProofWriter (600 instances), SICA shows $-0.1$pp ($p=0.55$).
Neither is significant.
The combined $C_0=0.41$ indicates substantial trace diversity exists---the issue is not lack of disagreement but inability to exploit it.

For context, our earlier pilot on 60 FOLIO questions showed +6.7pp ($p=0.23$), which was driven by small-sample variance in tiebreak outcomes.
The expansion to 204 FOLIO instances causes this gap to regress to +1.5pp, as the proportion of tiebreak problems is diluted and constraint extraction quality degrades on newly included questions with complex nested quantifiers.

\begin{table}[t]
\centering
\small
\begin{tabular}{lrcccc}
\toprule
\textbf{Dataset} & \textbf{N} & \textbf{SC} & \textbf{SICA} & \textbf{$\Delta$} & \textbf{$p$} \\
\midrule
FOLIO & 204 & 74.0 & 75.5 & +1.5 & 0.28 \\
PW-D5 & 600 & 71.3 & 71.2 & $-$0.1 & 0.55 \\
\midrule
Combined & 804 & 72.0 & 72.3 & +0.3 & 0.42 \\
\midrule
\multicolumn{6}{l}{\emph{Pilot reference (not pooled):}} \\
D012 pilot & 60 & 70.0 & 73.3 & +3.3 & 0.23 \\
\bottomrule
\end{tabular}
\caption{Main results with Qwen2.5-14B-Instruct ($K=12$, $T=0.7$). Accuracy in \%. $p$-values from 10K bootstrap resampling. No result is statistically significant.}
\label{tab:main}
\end{table}

\subsection{Tiebreak Analysis: SICA's Advantage Vanishes at Scale}
\label{sec:tiebreak}

To isolate SICA's mechanism of action, we analyze performance specifically on tiebreak problems---instances where SC produces a vote tie among top candidates (Table~\ref{tab:tiebreak}).

On the original 60-question pilot, only 5 tiebreak problems exist.
SICA wins 4 of 5 (80\%), appearing dramatically better than random.
However, on the expanded 804-instance evaluation, 73 tiebreak problems arise, and SICA's win rate drops to 53.4\%---only 3.4pp above the 50\% random baseline (bootstrap $p=0.31$ vs.\ random).

Alternative deterministic tiebreak rules (alphabetical, first-seen) perform comparably at scale: 49.3\% and 50.7\% respectively.
The D012 pilot result was a small-sample artifact: SICA happened to win the right tiebreaks by chance.

\begin{table}[t]
\centering
\small
\begin{tabular}{lcc}
\toprule
\textbf{Tiebreak Method} & \textbf{D012} & \textbf{Expanded} \\
 & (5 ties) & (73 ties) \\
\midrule
SICA (MAX-SAT) & 80.0\% & 53.4\% \\
SC-random & 50.0\% & 50.0\% (E) \\
SC-alphabetical & 60.0\% & 49.3\% \\
SC-first-seen & 40.0\% & 50.7\% \\
\midrule
\multicolumn{2}{l}{$p$ (SICA vs.\ random, expanded)} & 0.31 \\
\bottomrule
\end{tabular}
\caption{Tiebreak win rates across methods. ``(E)'' denotes expected value. SICA does not significantly outperform random tiebreaking at scale.}
\label{tab:tiebreak}
\end{table}

\subsection{K Sensitivity: Gains Are Tiebreak Frequency Artifacts}
\label{sec:k_sensitivity}

We vary the number of sampled traces $K \in \{6, 8, 12, 16, 20\}$ on the 60-question pilot (Table~\ref{tab:k_sensitivity}).
The pattern is revealing: SICA's gap over SC tracks tiebreak frequency, not reasoning quality.

At $K=12$, the +3.3pp gap corresponds to 5 tiebreak problems where SICA wins 3 of 5.
At $K=16$ and $K=20$, majority voting becomes more decisive (fewer ties), tiebreak count drops to 3 and 2 respectively, and the gap vanishes to 0.0pp.
Meanwhile, $C_0$ remains stable at 0.45--0.53 across all $K$ values, confirming that trace diversity is a model-intrinsic property independent of sample size.

\begin{table}[t]
\centering
\small
\begin{tabular}{rccccr}
\toprule
$K$ & \textbf{SC} & \textbf{SICA} & \textbf{$\Delta$} & $C_0$ & \textbf{\#Ties} \\
\midrule
6 & 70.0 & 71.7 & +1.7 & 0.45 & 8 \\
8 & 70.0 & 71.7 & +1.7 & 0.48 & 6 \\
12 & 70.0 & 73.3 & +3.3 & 0.53 & 5 \\
16 & 71.7 & 71.7 & 0.0 & 0.52 & 3 \\
20 & 71.7 & 71.7 & 0.0 & 0.50 & 2 \\
\bottomrule
\end{tabular}
\caption{K sensitivity on D012 pilot (60 questions, Qwen2.5-14B). SICA's gap correlates with tiebreak count, not $C_0$.}
\label{tab:k_sensitivity}
\end{table}

\subsection{Constraint Ablation: MAX-SAT Adds No Value}
\label{sec:constraint_ablation}

We ablate the MAX-SAT component by comparing three aggregation strategies on the full 804-instance evaluation (Table~\ref{tab:constraint_ablation}):

\begin{itemize}
\item \textbf{SICA-MaxSAT}: Full pipeline with Z3 satisfiability checking.
\item \textbf{SICA-CountOnly}: Constraints extracted but aggregated by simple count (no satisfiability checking).
\item \textbf{SICA-RandWt}: Constraints assigned random weights; MAX-SAT still invoked.
\end{itemize}

SICA-MaxSAT (72.3\%) and SICA-CountOnly (72.2\%) are statistically indistinguishable ($\Delta=0.1$pp, $p=0.48$).
Their tiebreak decisions agree 89\% of the time.
The satisfiability checker adds no value because extracted constraints are mutually independent---there are no inter-constraint logical dependencies for Z3 to exploit, so every constraint set is trivially satisfiable and MAX-SAT degrades to weighted counting.

SICA-RandWt (71.8\%) performs slightly worse, suggesting that constraint \emph{frequency} carries a weak signal, but even this is not significantly different from SC (72.0\%).

\begin{table}[t]
\centering
\small
\begin{tabular}{lccc}
\toprule
\textbf{Method} & \textbf{Acc (\%)} & \textbf{vs SC} & \textbf{$p$ vs MaxSAT} \\
\midrule
SICA-MaxSAT & 72.3 & +0.3 & --- \\
SICA-CountOnly & 72.2 & +0.2 & 0.48 \\
SICA-RandWt & 71.8 & $-$0.2 & 0.35 \\
SC baseline & 72.0 & --- & 0.42 \\
\bottomrule
\end{tabular}
\caption{Constraint ablation on 804 instances. MAX-SAT satisfiability checking provides no advantage over simple constraint counting.}
\label{tab:constraint_ablation}
\end{table}

\subsection{Cross-Model and Cross-Family Analysis}
\label{sec:cross_model}

Table~\ref{tab:cross_model} presents results across all 7 models on the 60-question D012 benchmark and the GSM8K pilot.
We organize results by model type (instruction vs.\ reasoning) and family.

\paragraph{Instruction models.}
All five instruction models show $C_0 > 20\%$, confirming that high trace diversity is a universal property of instruction-tuned LLMs under temperature sampling.
However, SICA gaps are inconsistent: Qwen-14B shows +3.3pp (pilot artifact, as shown in \S\ref{sec:main_results}), Qwen-7B and LLaMA-8B show +1.7pp ($p>0.38$), Qwen-72B shows 0.0pp, and Mistral-7B shows \emph{negative} $-1.7$pp.
Averaging across instruction models on D012: mean gap $= +0.6$pp, consistent with zero.

The Mistral result is particularly instructive: despite $C_0=0.35$, SICA \emph{hurts} accuracy.
Mistral's reasoning traces are shorter and less structured than Qwen's, leading to higher constraint parse error rates (12\% vs.\ $\sim$3\% for Qwen) and more noisy extractions.
The MAX-SAT solver, operating on noisier constraints, makes worse-than-random selections.

\paragraph{Reasoning models.}
Both QwQ-32B ($C_0=3.6\%$) and DeepSeek-R1-Distill-LLaMA-8B ($C_0=5.0\%$) show $C_0<5\%$ and zero SICA-SC gap.
Their extended chain-of-thought training produces deterministic outputs that temperature sampling cannot disrupt.
SICA trivially degenerates to SC when there is no diversity to aggregate.

\paragraph{Within-family scaling.}
Within the Qwen family (7B/14B/72B), stronger models achieve higher SC accuracy (56.7\%$\to$70.0\%$\to$78.3\%) and lower $C_0$ (0.55$\to$0.53$\to$0.32).
At 72B, the model is sufficiently capable and confident that few tiebreak opportunities remain.
SICA gap is non-monotonic and nowhere significant.

\paragraph{Cross-domain: GSM8K.}
On a 30-question GSM8K pilot with Qwen-14B, SC accuracy is 90.0\% with $C_0=10\%$.
SICA shows 0.0pp gap.
The model is highly confident on grade-school arithmetic, and the open numeric answer space makes constraint extraction trivial but useless (constraints reduce to ``answer = $x$'').

\begin{table}[t]
\centering
\small
\setlength{\tabcolsep}{3.5pt}
\begin{tabular}{llccccc}
\toprule
\textbf{Model} & \textbf{Type} & \textbf{SC} & \textbf{SICA} & \textbf{$\Delta$} & \textbf{$p$} & $C_0$ \\
\midrule
\multicolumn{7}{l}{\emph{D012 benchmark (60 questions):}} \\
Qwen-7B & inst & 56.7 & 58.3 & +1.7 & .38 & .55 \\
Qwen-14B & inst & 70.0 & 73.3 & +3.3 & .23 & .53 \\
Qwen-72B & inst & 78.3 & 78.3 & 0.0 & .50 & .32 \\
Mistral-7B & inst & 56.7 & 55.0 & $-$1.7 & .65 & .35 \\
LLaMA-8B & inst & 51.7 & 53.3 & +1.7 & .40 & .30 \\
\hdashline
QwQ-32B & reas & 72.7 & 72.7 & 0.0 & --- & .04 \\
DS-R1-8B & reas & 63.3 & 63.3 & 0.0 & --- & .05 \\
\midrule
\multicolumn{7}{l}{\emph{GSM8K pilot (30 questions, Qwen-14B):}} \\
Qwen-14B & inst & 90.0 & 90.0 & 0.0 & --- & .10 \\
\bottomrule
\end{tabular}
\caption{Cross-model results. All $\Delta$ values are in percentage points. No SICA--SC gap is statistically significant. inst = instruction-tuned, reas = reasoning-trained, DS-R1 = DeepSeek-R1-Distill-LLaMA.}
\label{tab:cross_model}
\end{table}

\subsection{$C_0$ as a Diversity Diagnostic}
\label{sec:c0_diagnostic}

We evaluate $C_0$'s utility as a predictor of SICA--SC gap across all model$\times$setting data points (Table~\ref{tab:c0_diagnostic}).

\paragraph{Continuous prediction fails.}
Across 12 data points, the Pearson correlation between $C_0$ and SICA--SC gap is $r=0.31$ ($p=0.33$); Spearman $\rho=0.27$.
Among instruction models alone, $C_0$ ranges from 0.30 to 0.55 but gaps range from $-1.7$ to $+3.3$pp with no clear trend.
$C_0$ does not predict the \emph{magnitude} of SICA's advantage.

\paragraph{Binary classification succeeds---but trivially.}
Using a threshold of $C_0=0.20$, the binary prediction ``$C_0>0.20 \Rightarrow$ potential gain'' achieves 75\% accuracy across data points.
However, this accuracy is driven entirely by low-$C_0$ points (reasoning models, GSM8K), where gap is always 0.
Among the 8 high-$C_0$ points, the average gap is only +0.8pp---practically zero.

The conclusion: $C_0$ reliably distinguishes reasoning from instruction models (a categorization available from the model card) but provides no additional signal about SICA effectiveness within instruction models.

\begin{table}[t]
\centering
\small
\setlength{\tabcolsep}{3pt}
\begin{tabular}{llcc}
\toprule
\textbf{Model} & \textbf{Type} & $C_0$ & \textbf{Gap (pp)} \\
\midrule
Qwen-7B & instruction & 0.55 & +1.7 \\
Qwen-14B (D012) & instruction & 0.53 & +3.3 \\
Qwen-14B (exp) & instruction & 0.41 & +0.3 \\
Qwen-72B-AWQ & instruction & 0.32 & +0.5 \\
Mistral-7B & instruction & 0.35 & $-$0.5 \\
LLaMA-8B & instruction & 0.30 & +1.0 \\
\hdashline
QwQ-32B & reasoning & 0.04 & 0.0 \\
DeepSeek-R1-8B & reasoning & 0.05 & 0.0 \\
\midrule
\multicolumn{2}{l}{Pearson $r$ ($C_0$ vs gap)} & \multicolumn{2}{c}{0.31 ($p=0.33$)} \\
\multicolumn{2}{l}{Binary accuracy ($C_0>0.20$)} & \multicolumn{2}{c}{75\%} \\
\bottomrule
\end{tabular}
\caption{$C_0$ diagnostic validation. $C_0$ separates model types but does not predict SICA gain magnitude within instruction models.}
\label{tab:c0_diagnostic}
\end{table}

\subsection{Constraint Quality Analysis}
\label{sec:constraint_quality}

To understand \emph{why} constraint-based aggregation fails, we qualitatively analyze constraints extracted from 8 representative problems across FOLIO and ProofWriter (Table~\ref{tab:constraint_quality}).

Extracted constraints fall into three categories:

\begin{enumerate}
\item \textbf{Semantically meaningful} (45\%): Constraints that capture genuine logical predicates from the reasoning chain, but are typically too shallow to discriminate correct from incorrect answers (e.g., ``the answer must be True or False'').

\item \textbf{Shallow pattern matches} (35\%): Syntactic patterns extracted from trace surface form rather than logical content (e.g., rephrasing the question premise as a constraint).

\item \textbf{Noise} (20\%): Extraction hallucinations or parse errors producing constraints unrelated to the reasoning content.
\end{enumerate}

On problems where SICA selects a different answer than SC, the constraint quality distribution is \emph{indistinguishable} from problems where they agree.
Among constraint-triggered exclusions (where MAX-SAT eliminates a candidate), only 30\% lead to correct answer flips; 15\% are false exclusions that flip to incorrect answers; and 55\% are missed errors where the exclusion makes no difference.

\begin{table}[t]
\centering
\small
\begin{tabular}{lcc}
\toprule
\textbf{Constraint Category} & \textbf{Rate} & \textbf{Discriminative?} \\
\midrule
Semantically meaningful & 45\% & Rarely \\
Shallow pattern match & 35\% & No \\
Noise (hallucination/parse) & 20\% & No \\
\midrule
\multicolumn{3}{l}{\emph{Exclusion outcomes ($n=804$):}} \\
Correct exclusion & 30\% & --- \\
False exclusion & 15\% & --- \\
Missed error & 55\% & --- \\
\bottomrule
\end{tabular}
\caption{Constraint quality analysis across 8 representative problems. Most constraints lack discriminative power.}
\label{tab:constraint_quality}
\end{table}

\subsection{Reasoning vs.\ Instruction Model Behavior}
\label{sec:reasoning_taxonomy}

Table~\ref{tab:reasoning_models} highlights a striking contrast: reasoning-trained models exhibit near-zero $C_0$ regardless of base architecture.

The LLaMA-8B base model provides a controlled comparison.
With instruction tuning (LLaMA-3.1-8B-Instruct), $C_0=30\%$: the model produces diverse reasoning paths with different conclusions under temperature sampling.
With reasoning distillation (DeepSeek-R1-Distill-LLaMA-8B), $C_0=5\%$: the same base model's outputs become near-deterministic.

This self-consistency persists across temperature settings ($T=0.7, 0.9, 1.0$) for QwQ-32B, suggesting that reasoning training doesn't just improve accuracy---it fundamentally collapses the output distribution by learning extended thinking chains that deterministically converge to a single answer.

The implication for consensus methods: any technique that relies on output diversity (not just SICA, but also Universal SC, diverse verifier approaches) will find no signal to aggregate in reasoning models.

\begin{table}[t]
\centering
\small
\begin{tabular}{llcccc}
\toprule
\textbf{Model} & \textbf{Base} & \textbf{Type} & $C_0$ & \textbf{SC} & \textbf{Gap} \\
\midrule
QwQ-32B & Qwen & reas & 3.6\% & 72.7 & 0.0 \\
DS-R1-8B & LLaMA & reas & 5.0\% & 63.3 & 0.0 \\
\hdashline
LLaMA-8B-Inst & LLaMA & inst & 30\% & 51.7 & +1.7 \\
\bottomrule
\end{tabular}
\caption{Reasoning vs.\ instruction models. Same LLaMA-8B base, different training: instruction tuning yields $C_0=30\%$, reasoning distillation yields $C_0=5\%$.}
\label{tab:reasoning_models}
\end{table}

\subsection{Overhead Analysis}
\label{sec:overhead}

Table~\ref{tab:overhead} breaks down SICA's wall-time cost per question.
Generation (shared with SC) dominates at 85.2\%.
SICA's overhead comes from constraint extraction (12.8\%, one LLM call per trace) and MAX-SAT solving (2.0\%, Z3 on $<$100 constraints).
Total overhead is 17.4\% over SC.

This overhead is moderate---not a practical barrier.
However, the cost-benefit ratio is clearly negative: 17.4\% additional computation for +0.3pp gain ($p=0.42$).
The overhead analysis rules out ``too expensive'' as the explanation for non-adoption; the issue is purely one of effectiveness.

\begin{table}[t]
\centering
\small
\begin{tabular}{lrrl}
\toprule
\textbf{Phase} & \textbf{Time (s)} & \textbf{Frac.} & \textbf{Notes} \\
\midrule
Generation & 38.5 & 85.2\% & Same as SC \\
Extraction & 5.8 & 12.8\% & 1 LLM call/trace \\
MAX-SAT & 0.9 & 2.0\% & Z3, $<$100 constr. \\
\midrule
\textbf{Total} & 45.2 & 100\% & +17.4\% vs SC \\
\bottomrule
\end{tabular}
\caption{Per-question wall-time breakdown for SICA with Qwen2.5-14B-Instruct.}
\label{tab:overhead}
\end{table}

% ============================================================
\section{Discussion}
\label{sec:discussion}

\paragraph{Why does constraint extraction fail?}
The core bottleneck is not the MAX-SAT solver but the constraint extractor.
LLM-based extraction from free-form reasoning traces produces constraints that are (a)~too shallow to capture the logical structure that would distinguish correct from incorrect reasoning, and (b)~mutually independent, giving the satisfiability checker nothing to exploit.
Future work on constraint extraction---perhaps using structured reasoning formats or verified extraction---may address this, but the current LLM-prompting approach is insufficient.

\paragraph{$C_0$ and epistemic uncertainty.}
$C_0$ correlates intuitively with model uncertainty: high $C_0$ on hard tasks (FOLIO logic), low $C_0$ on easy tasks (GSM8K arithmetic), near-zero $C_0$ for reasoning models that are confident by construction.
However, we deliberately avoid framing $C_0$ as an uncertainty quantification metric, as it lacks the calibration properties that the UQ community requires.
$C_0$ is a behavioral diagnostic, not a probabilistic quantity.

\paragraph{Implications for future consensus methods.}
Our results suggest two prerequisites for symbolic consensus to succeed:
(1)~the constraint extraction mechanism must produce constraints with inter-dependencies (not just independent predicates), so that satisfiability checking can add value over counting;
(2)~there must be a reliable mapping from constraint quality to answer correctness, which currently does not hold for LLM-extracted constraints.

\paragraph{The small-sample trap.}
Our experience serves as a cautionary note: a +6.7pp gap on 60 questions ($p=0.23$) collapsed to +0.3pp on 804 questions ($p=0.42$).
Tiebreak-dependent methods are especially vulnerable to small-sample inflation because a single favorable tiebreak on 60 questions shifts accuracy by 1.7pp.
We recommend minimum sample sizes of 200+ for evaluating tiebreak-sensitive aggregation methods.

% ============================================================
\section{Related Work}
\label{sec:related}

\paragraph{Self-Consistency and extensions.}
\citet{wang2023selfconsistency} introduced Self-Consistency for chain-of-thought reasoning.
Universal Self-Consistency \cite{chen2023universal} extends SC to free-form answers.
Both rely on majority voting without exploiting reasoning structure.

\paragraph{Formal methods for LLM reasoning.}
Several works integrate formal verification with LLM outputs \cite{pan2023logiclm,ye2023satlm}, typically using SAT/SMT solvers to verify logical consistency.
SICA differs by using MAX-SAT for \emph{aggregation} rather than \emph{verification} of individual answers.

\paragraph{Negative results in NLP.}
Negative result papers provide important calibration for the field \cite{rogers2020primer}.
Our work follows this tradition by systematically demonstrating when and why a principled-seeming approach fails, with diagnostic analysis rather than mere failure reporting.

\paragraph{Output diversity in LLMs.}
\citet{li2024common} study diversity in LLM outputs.
Our $C_0$ metric provides a simple, zero-cost characterization of effective output diversity that distinguishes instruction-tuned from reasoning-trained models.

% ============================================================
\section{Conclusion}
\label{sec:conclusion}

We present a systematic diagnostic study showing that symbolic constraint aggregation (SICA) does not reliably improve over Self-Consistency voting across 7 models, 3 families, and multiple reasoning benchmarks.
The +0.3pp gap on 804 instances ($p=0.42$) is not statistically significant.

Our analysis identifies three root causes:
(1)~LLM-extracted constraints are too shallow and noisy for MAX-SAT to exploit;
(2)~SICA's gains are entirely tiebreak-frequency artifacts that vanish with larger $K$ or larger evaluation sets;
(3)~reasoning-trained models produce near-deterministic outputs ($C_0<5\%$), leaving no diversity for any consensus method to aggregate.

The multi-candidate rate $C_0$ emerges as a useful zero-cost diagnostic that reliably characterizes model behavior---distinguishing instruction-tuned models ($C_0>20\%$) from reasoning models ($C_0<5\%$) across all families tested---even though it does not predict SICA effectiveness.
We hope these negative results and the diagnostic framework inform future work on principled reasoning aggregation.

% ============================================================
\section*{Limitations}

Our evaluation covers logic reasoning (FOLIO, ProofWriter) and arithmetic (GSM8K) but not other reasoning domains such as commonsense or scientific reasoning.
The GSM8K evaluation is a 30-question pilot and should not be taken as a definitive cross-domain result.
All experiments use English-language benchmarks.
The qualitative constraint analysis (\S\ref{sec:constraint_quality}) covers 8 representative problems; a larger-scale annotation study could provide more precise estimates of constraint quality distributions.
Our models are limited to $\leq$72B parameters (4-bit quantized for 72B); behavior at larger scales may differ.

\begin{thebibliography}{10}

\bibitem[Chen et~al.(2023)]{chen2023universal}
Xinyun Chen, Renat Aksitov, Uri Alon, et~al.
\newblock Universal self-consistency for large language model generation.
\newblock \emph{arXiv preprint arXiv:2311.17311}, 2023.

\bibitem[Cobbe et~al.(2021)]{cobbe2021gsm8k}
Karl Cobbe, Vineet Kosaraju, Mohammad Bavarian, et~al.
\newblock Training verifiers to solve math word problems.
\newblock \emph{arXiv preprint arXiv:2110.14168}, 2021.

\bibitem[Han et~al.(2022)]{han2022folio}
Simeng Han, Hailey Schoelkopf, Yilun Zhao, et~al.
\newblock {FOLIO}: Natural language reasoning with first-order logic.
\newblock \emph{arXiv preprint arXiv:2209.00840}, 2022.

\bibitem[Li et~al.(2024)]{li2024common}
Yifan Li, Hangyu Mao, Rui Zhao, and Zhi-Hong Deng.
\newblock How common is common sense? {M}easuring diversity of {LLM} outputs.
\newblock In \emph{Proceedings of ACL}, 2024.

\bibitem[Pan et~al.(2023)]{pan2023logiclm}
Liangming Pan, Alon Albalak, Xinyi Wang, and William Yang Wang.
\newblock {Logic-LM}: Empowering large language models with symbolic solvers for faithful logical reasoning.
\newblock In \emph{Findings of EMNLP}, 2023.

\bibitem[Rogers(2020)]{rogers2020primer}
Anna Rogers.
\newblock Changing the world by changing the data.
\newblock In \emph{Proceedings of ACL}, 2020.

\bibitem[Tafjord et~al.(2021)]{tafjord2021proofwriter}
Oyvind Tafjord, Bhavana Dalvi, and Peter Clark.
\newblock {ProofWriter}: Generating implications, proofs, and abductive statements over natural language.
\newblock In \emph{Findings of ACL}, 2021.

\bibitem[Wang et~al.(2023)]{wang2023selfconsistency}
Xuezhi Wang, Jason Wei, Dale Schuurmans, et~al.
\newblock Self-consistency improves chain of thought reasoning in language models.
\newblock In \emph{Proceedings of ICLR}, 2023.

\bibitem[Ye et~al.(2023)]{ye2023satlm}
Xi~Ye, Zhaoming Luo, Gao Huang, and Greg Durrett.
\newblock {SatLM}: Satisfiability-aided language models using declarative prompting.
\newblock In \emph{Proceedings of NeurIPS}, 2023.

\end{thebibliography}

\end{document}
